Zeichnen Sie die Wege mit den Parametrisierungen mit $t\in I=[0,1]$
\begin{teilaufgaben}
\item $\gamma_1(t) = 1 + it$
\item $\gamma_2(t) = e^{-\pi it}$
\item $\gamma_3(t) = e^{\pi it}$
\item $\gamma_4(t) = 1 + it + t^2$
\end{teilaufgaben}
Berechnen Sie ausserdem die Ableitung nach $t$.

\begin{loesung}
\begin{figure}
\centering
\begin{tikzpicture}[>=latex,thick,scale=3]
\def\w{70}

\draw[->] (-1.1,0) -- (2.1,0) coordinate[label={$\operatorname{Re}$}];
\draw[->] (0,-1.1) -- (0,1.2) coordinate[label={right:$\operatorname{Im}$}];

\draw[->,color=darkred,line width=1.2pt] (1,0) -- (1,1);
\node at (1,0.7) [right] {$\gamma_1$};

\draw[color=darkred,line width=1.2pt] (1,0) arc(0:-179:1);
\draw[<-,color=darkred,line width=1.2pt] (-1,0) -- ++(0,-0.03);
\node at (-\w:1) [below right] {$\gamma_2$};

\draw[color=darkred,line width=1.2pt] (1,0) arc(0:179:1);
\draw[<-,color=darkred,line width=1.2pt] (-1,0) -- ++(0,0.03);
\node at (\w:1) [above right] {$\gamma_3$};

\draw[->,color=darkred,line width=1.2pt]
	plot[domain=0:1,samples=100] ({1+\x*\x},{\x});
\def\X{0.8}
\node at ({1+\X*\X},{\X}) [below right] {$\gamma_4$};

\fill[color=darkred] (1,0) circle[radius=0.03];

\end{tikzpicture}
\caption{Pfade von Aufgabe \ref{301}.
\label{buch:integration:aufgabe:fig:301}}
\end{figure}
Die Pfade sind in Abbildung~\ref{buch:integration:aufgabe:fig:301}
dargestellt.
Die Ableitungen sind die folgenden:
\begin{teilaufgaben}
\item $\dot{\gamma}_1(t) = i$.
\item $\dot{\gamma}_2(t) = -\pi i e^{-\pi it}$.
\item $\dot{\gamma}_3(t) =  \pi i e^{\pi it}$.
\item $\dot{\gamma}_4(t) = 2t+it$.
\qedhere
\end{teilaufgaben}
\end{loesung}
